% =====================================================================
%  統計基礎 3. ベイズ理論
%  構成は Docs/SectionPlan.md にしたがう。
% =====================================================================
\chapter{ベイズ理論}
\label{chap:stat-bayes}

\begin{chaptergoal}
  % TODO: この章の目標
\end{chaptergoal}

% ---------------------------------------------------------------------
\section{事前分布・事後分布}
\label{sec:stat-prior-posterior}

\keywords{事前分布、共役事前分布、事後分布}

% TODO: 〔補足〕ベイズ統計の考え方
% TODO:   ベイズの定理の復習
% TODO:   パラメータを確率変数として扱う。これまでの（頻度論の）考え方との違い

\subsection{事前分布}
% TODO:   データを見る前の知識や信念を分布で表す
% TODO:   情報が少ないときの事前分布（無情報事前分布）

\subsection{事後分布}
% TODO:   事後分布 ∝ 尤度 × 事前分布
% TODO:   正規化定数（全体の確率を 1 にする定数）の役割
% TODO:   データが増えると事前分布の影響が小さくなること

\subsection{共役事前分布}
% TODO:   定義：事前分布と事後分布が同じ種類の分布になる組み合わせ
% TODO:   ベータ分布と二項分布（コイン投げの例）
% TODO:   ガンマ分布とポアソン分布
% TODO:   正規分布と正規分布（分散が既知の場合）
% TODO: 〔補足〕事後分布のまとめ方：事後平均、事後確率最大の値（MAP 推定）、信用区間（信頼区間との違い）

% ---------------------------------------------------------------------
\section{ベイズ的仮説検定}
\label{sec:stat-bayes-testing}

\keywords{ベイズファクター、ベイズ判別（各カテゴリーの事後確率）}

% TODO: 考え方：仮説そのものが正しい確率（事後確率）を比べる
% TODO: 〔補足〕周辺尤度（パラメータについて平均した尤度）

\subsection{ベイズファクター}
% TODO:   定義：2つの仮説（モデル）の周辺尤度の比
% TODO:   事後オッズ ＝ ベイズファクター × 事前オッズ
% TODO:   値の大きさの目安（ジェフリーズの基準）
% TODO:   p 値との違い

\subsection{ベイズ判別}
% TODO:   各カテゴリーの事後確率を、ベイズの定理で計算する
% TODO:   事後確率が最大のカテゴリーに分類する（誤分類の確率が最小になる）
% TODO:   例：2つの正規分布のどちらから来たデータかを判別する
% TODO:   モデリング側の「単純ベイズ」「ベイズ判別」への橋渡し

\begin{chaptersummary}
  % TODO: この章のまとめ
\end{chaptersummary}
